\section*{Chapter \ref{ch:g12:contdist} --- Continuous Random Variables}

\begin{solution}{exo:g12:contdist:1}
$\P(X \leq 5) = \frac{5}{15} = \frac13$;
$\P(4 \leq X \leq 10) = \frac{6}{15} = \frac25$;
$\E(X) = \frac{15}{2} = 7.5$ min;
$\sigma(X) = \frac{15}{\sqrt{12}} = \frac{5\sqrt3}{2} \approx 4.3$ min.
\end{solution}

\begin{solution}{exo:g12:contdist:2}
$f \geq 0$ and
$\int_0^2 \frac38 t^2\,\dd t = \frac38\left[\frac{t^3}{3}\right]_0^2
= \frac38 \times \frac83 = 1$: a \omterm{def:g12:contdist:density}{density}.
\[
\E(X) = \int_0^2 \frac38 t^3\,\dd t = \frac38 \times \frac{16}{4} = \frac32,
\qquad
\P(X \geq 1) = \int_1^2 \frac38 t^2\,\dd t = \frac{8 - 1}{8} = \frac78 .
\]
\end{solution}

\begin{solution}{exo:g12:contdist:3}
\emph{1.} $\E(X) = \frac{1}{0.2} = 5$ years;
$\P(X > 5) = \eu^{-0.2\times5} = \eu^{-1} \approx 0.37$.

\emph{2.} By memorylessness (\cref{thm:g12:contdist:memoryless}),
$\pcond{X > 3}{X > 8} = \P(X > 5) = \eu^{-1} \approx 0.37$: the three years
of service change nothing.
\end{solution}

\begin{solution}{exo:g12:contdist:4}
$\P(X > m) = \eu^{-\lambda m} = \frac12$ gives
$m = \frac{\ln 2}{\lambda} \approx \frac{0.69}{\lambda}$, smaller than
$\E(X) = \frac1\lambda$. The exponential \omterm{def:g12:contdist:density}{density} has a long right tail: a
few atypically long lifetimes pull the \emph{\omterm{def:g11:stat:mean}{mean}} above the
\emph{\omterm{def:g11:stat:median}{median}}, the value that half the population exceeds. (This is the
half-life of \cref{ex:g12:exp:halflife}: half the atoms survive it, even
though the average lifetime is longer.)
\end{solution}

\begin{solution}{exo:g12:contdist:5}
$168 = \mu - \sigma$ and $182 = \mu + \sigma$: about $68\%$.
$189 = \mu + 2\sigma$: above it lies half of the remaining
$100 - 95.4 = 4.6\%$, so about $2.3\%$.
$154 = \mu - 3\sigma$: about $\frac{100 - 99.7}{2} = 0.15\%$.
\end{solution}

\begin{solution}{exo:g12:contdist:6}
$\P(X < 500) \leq 2.5\%$ means $500$ must sit at least two \omterm{def:g11:stat:variance}{standard
deviations} below the \omterm{def:g11:stat:mean}{mean} (the normal leaves $2.3\% \approx 2.5\%$ below
$\mu - 2\sigma$; with the conventional $1.96$, the reasoning is identical):
\[
\mu - 2\sigma \geq 500 \iff \mu \geq 500 + 8 = 508 \text{ g}.
\]
The machine must be set to $508$ g on average --- the price of the
guarantee is $8$ g of ``free'' product per bag.
\end{solution}

\begin{solution}{exo:g12:contdist:7}
\emph{1.} With $p = \frac16$, $n = 1200$:
$\sqrt{\frac{p(1-p)}{n}} = \sqrt{\frac{5/36}{1200}} \approx 0.0108$, so the
$95\%$ fluctuation \omterm{def:g10:numbers:interval}{interval} is
\[
\frac16 \pm 1.96 \times 0.0108 \approx \intcc{0.146}{0.188}.
\]

\emph{2.} The observed \omterm{def:g10:stats:series}{frequency} $\frac{260}{1200} \approx 0.217$ lies far
outside: the fairness hypothesis is rejected at the $5\%$ level. The
deviation is about $4.6$ \omterm{def:g11:stat:variance}{standard deviations} --- for a normal
\omterm{def:g10:numbers:approx}{approximation}, a \omterm{def:g10:proba:distribution}{probability} of order $10^{-6}$, far more conclusive than
the $\leq 4.6\%$ bound from Bienaymé--Chebyshev
(\cref{exo:g12:sums:7}).
\end{solution}

\begin{solution}{exo:g12:contdist:8}
\emph{1.} Margin $\frac{1}{\sqrt n} \leq 0.02$ requires
$n \geq 2500$ people.

\emph{2.} To distinguish scores $1$ point apart, the margin must be well
under $0.5$ point, requiring $n \geq \left(\frac{1}{0.005}\right)^2 =
40\,000$ by the simplified formula --- already impractical for most polls.
Worse, the statistical margin only accounts for \emph{sampling} error;
systematic biases (unrepresentative \omterm{def:g10:proba:sample}{samples}, non-response, last-minute
swings) do not shrink as $n$ grows and typically exceed one point. A
$1$-point race is genuinely too close to call.
\end{solution}

\begin{solution}{exo:g12:contdist:9}
\emph{1.} Since $\alpha > 0$, $c \leq \alpha X + \beta \leq d
\iff \frac{c - \beta}{\alpha} \leq X \leq \frac{d - \beta}{\alpha}$, so
\[
\P(c \leq Y \leq d)
= \int_{(c-\beta)/\alpha}^{(d-\beta)/\alpha} f(t)\,\dd t .
\]
The substitution $y = \alpha t + \beta$ (i.e.\ reading the area in the $y$
variable, $t = \frac{y - \beta}{\alpha}$, $\dd t = \frac{\dd y}{\alpha}$)
turns this into
$\int_c^d \frac{1}{\alpha} f\left(\frac{y-\beta}{\alpha}\right)\dd y$: the
\omterm{def:g11:func:function}{function} $g(y) = \frac1\alpha f\left(\frac{y-\beta}{\alpha}\right)$ is a
\omterm{def:g12:contdist:density}{density} of $Y$.

\emph{2.} With $f = \varphi$, $\alpha = \sigma$, $\beta = \mu$:
\[
g(y) = \frac{1}{\sigma}\,\varphi\!\left(\frac{y - \mu}{\sigma}\right)
= \frac{1}{\sigma\sqrt{2\pi}}\,
\exp\left(-\frac{(y-\mu)^2}{2\sigma^2}\right),
\]
which is therefore the \omterm{def:g12:contdist:density}{density} of
$\mathcal N(\mu, \sigma^2) = \mu + \sigma\,\mathcal N(0,1)$.
\end{solution}

\begin{solution}{pb:g12:contdist:1}
\textbf{1.} $\P(\text{wait} \leq 10) = \frac{10}{30} =
\frac13$; $\E = 15$ min;
$\sigma = \frac{30}{\sqrt{12}} \approx 8.7$ min.

\textbf{2.} $\int_0^1 2x\,\dd x = 1$: a \omterm{def:g12:contdist:density}{density}.
$\E(X) = \int_0^1 2x^2\,\dd x = \frac23$;
$\P\left(X > \frac12\right) = \int_{1/2}^1 2x\,\dd x =
1 - \frac14 = \frac34$.

\textbf{3.} $\P(T > 15) = \eu^{-15/10} \approx 0.22$. \omterm{def:g11:stat:median}{Median}:
$\eu^{-m/10} = \frac12$: $m = 10\ln 2 \approx 6.9$ min ---
less than the \omterm{def:g11:stat:mean}{mean} $10$ because the exponential's long right
tail (rare huge waits) drags the \omterm{def:g11:stat:mean}{mean} upward while leaving the
\omterm{def:g11:stat:median}{median} at the ``typical'' wait.

\textbf{4.} By memorylessness,
$\P(T > 35 \mid T > 20) = \P(T > 15) \approx 0.22$: the
twenty minutes already served buy nothing --- the stream does
not age.

\textbf{5.} (a) exponential (radioactive decay is the model
memorylessness was invented for); (b) uniform on
$\intcc{0}{8}$ (a schedule with a random offset); (c)
exponential; (d) neither: a worn bulb is \emph{more} likely to
fail in the next hour than a new one, so
$\P(T > s + t \mid T > s) < \P(T > t)$: aging contradicts the
memorylessness theorem.

\textbf{6.} $68\,\%$, $95\,\%$, $99.7\,\%$: the three
landmark \omterm{def:g10:numbers:interval}{intervals}.

\textbf{7.} $z = \frac{190 - 172}{8} = 2.25$: upper tail
$\approx 1.2\,\%$.

\textbf{8.} $z = 2$: about $2.3\,\%$ --- roughly one person
in $44$.

\textbf{9.} The spec sits at $\pm 2\sigma$: about
$95.4\,\%$ pass, $4.6\,\%$ rejected --- ruinous at scale. At
$\pm 6\sigma$ the failure rate drops to about two parts per
\emph{billion}: six-sigma buys the right to mass-produce
without mass-inspecting.

\textbf{10.} $\sigma = 5$; with \omterm{def:g12:limcont:continuity}{continuity} correction,
$\P \approx \P\left(\abs Z \leq \frac{5.5}{5}\right) =
\P(\abs Z \leq 1.1) \approx 0.73$. The theorem smooths the
Galton board of \cref{pb:g11:binom:1}: the staircase of slots
becomes the \omterm{def:g12:limcont:continuity}{continuous} bell.

\textbf{11.} $n \approx \frac{1.96^2 \times 0.25}{0.005^2}
\approx 38\,400$ people --- and at that size the
\emph{non}-sampling errors (frame, refusals, lies) dwarf the
margin: a one-point race is beyond honest polling, which is
why serious institutes then say ``too close to call''.

\textbf{12.} Memorylessness: from the moment you arrive, the
remaining wait is exponential with \omterm{def:g11:stat:mean}{mean} $10$ --- the full
$10$ minutes, not $5$. The timetable's ``\omterm{def:g11:stat:mean}{mean} gap $10$''
misleads: your wait has the same \omterm{def:g12:randvar:rv}{distribution} as a whole gap.

\textbf{13.} Random arrivals land proportionally to gap
length: $\P(\text{long gap}) = \frac{15}{20} = \frac34$.
Expected wait: $\frac34 \times \frac{15}{2} + \frac14 \times
\frac52 = 5.625 + 0.625 = 6.25$ min --- more than the naive
$5$, though the \omterm{def:g11:stat:mean}{mean} gap is $10$. Culprit: \emph{length-biased
sampling} --- long gaps catch more passengers.

\textbf{14.} Backward and forward waits are each exponential
with \omterm{def:g11:stat:mean}{mean} $10$ (memorylessness runs in both directions from
your arrival): the gap containing you averages $20$ minutes
--- twice the typical gap. Inspection paradox in one sentence:
\emph{the \omterm{def:g10:numbers:interval}{interval} you happen to inspect is not a typical
\omterm{def:g10:numbers:interval}{interval}, because you were more likely to fall into a big
one.}

\textbf{15.} Average class: $\frac{9 \times 10 + 110}{10} =
20$ students. Student-experienced average:
$\frac{90 \times 10 + 110 \times 110}{200} = \frac{13\,000}
{200} = 65$ students. The brochure prints $20$; three in five
students sit in the class of $110$ and live the $65$.

\textbf{16.} Popular people appear on many friend-lists, so
sampling ``a friend'' is length-biased towards the sociable
--- your friends are drawn from the long gaps of the social
timetable.

\textbf{17.} For $t \geq 0$:
$\P(T > t) = \P(-10\ln U > t) = \P\left(U <
\eu^{-t/10}\right) = \eu^{-t/10}$: exactly the exponential's
survival \omterm{def:g11:func:function}{function} --- one \omterm{def:g12:exp:ln}{logarithm} converts the computer's
uniform noise into any waiting time.

\textbf{18.} Each uniform has \omterm{def:g11:stat:mean}{mean} $\frac12$ and \omterm{def:g12:randvar:exp}{variance}
$\frac{1}{12}$: the sum of twelve has \omterm{def:g11:stat:mean}{mean} $6$ and \omterm{def:g12:randvar:exp}{variance}
$1$ --- so the recipe outputs \omterm{def:g11:stat:mean}{mean} $0$, \omterm{def:g12:randvar:exp}{variance} $1$. Twelve
\omterm{def:g12:sums:indep}{independent} small pushes summed: the Galton mechanism, and by
De Moivre--Laplace's blessing the histogram is already
bell-shaped to the eye.

\textbf{19.} About the model: markets are not sums of many
\emph{\omterm{def:g12:sums:indep}{independent}} small causes --- panics correlate
everything (the 2008 lesson of the previous problem) and
produce fat tails no normal \omterm{def:g12:contdist:density}{density} owns. When the data
whispers $10^{-137}$, the statistician's catechism answers:
\emph{the model is rejected, not the day}.

\textbf{20.} Uniform: all positions equal, the honest ``I
know only the bounds''. Exponential: the hazard that never
ages --- decays, arrivals, clicks. Normal: the democratic
bell of many small \omterm{def:g12:sums:indep}{independent} causes --- heights, errors,
averages. Sting: what you \omterm{def:g10:proba:sample}{sample} is biased by how you bumped
into it --- buses, classes, friends. And the arc: the child
who counted juice cartons in grade 6 has, ten volumes of
problems later, measured the crowd with a curve; the epsilon,
the \omterm{def:g12:integ:area}{integral} and the central limit theorem are waiting in the
university volumes to explain \emph{why} the bell tolls for
everything.
\end{solution}
